\section{Trace data constrains restricted sections}

A restricted section must vanish along the restricted multiplier ideal. Passing to the asymptotic orders of \cref{chap:trace-data} gives a degree obstruction on every smooth compact modification of the submanifold. Subtracting any finite collection of these trace orders also leaves a pseudoeffective class. Here $c_1$ denotes the real Bott--Chern first Chern class.

Let $Y\subset X$ be smooth, compact and connected, of positive dimension $r$, with $\nu(u,Y)=0$. For a positive integer $m$, write
\[
 \mathfrak a_m=\mathcal I(mu)\mathcal O_Y,
 \qquad
 V_m=\operatorname{im}\bigl(\mathcal H_m(u)\longrightarrow H^0(Y,L^m|_Y)\bigr).
\]
Thus $V_m\subseteq H^0(Y,L^m|_Y\otimes\mathfrak a_m)$; the right-hand space can contain sections that do not extend to $X$.

\Needspace{18\baselineskip}
\begin{theorem}\label{thm:gauduchon-trace-section-obstruction}
Suppose $V_m\neq0$ for arbitrarily large integers $m$. Let $\pi:Z\to Y$ be any smooth compact modification, let $E_1,\ldots,E_q$ be distinct prime divisors on $Z$, and let $\gamma$ be a Gauduchon metric on $Z$. Then
\[
 \deg_\gamma(\pi^*(L|_Y))
 \geq\sum_{i=1}^q\tau_{Y,u}(E_i)\int_{E_i}\gamma^{r-1}.
\]
Moreover the real closed $(1,1)$ class
\[
 c_1(\pi^*(L|_Y))-
       \sum_{i=1}^q\tau_{Y,u}(E_i)c_1(\mathcal O_Z(E_i))
\]
is pseudoeffective: it admits a closed positive current.
\end{theorem}
\begin{proof}
Choose a nonzero $s_m\in V_m$ for each integer in an unbounded sequence of such degrees. Its pullback is a nonzero holomorphic section of $\pi^*(L^m|_Y)$. Set $b_i(m)=\operatorname{ord}_{E_i}\mathfrak a_m$, the generic order of the pulled-back image ideal along $E_i$. Since the local coefficient of the pulled-back section belongs to $\mathfrak a_m$, its zero divisor satisfies
\[
 \operatorname{div}(\pi^*s_m)
 \geq\sum_{i=1}^q b_i(m)E_i.
\]
Poincar\'e--Lelong and the Gauduchon identity, exactly as in \cref{thm:gauduchon-section-obstruction}, give
\[
 m\deg_\gamma(\pi^*(L|_Y))
 \geq\sum_{i=1}^q b_i(m)\int_{E_i}\gamma^{r-1}.
\]
Divide by $m$ and use the full integer limit in \cref{thm:intrinsic-trace-profile}. The same limit holds on the selected unbounded sequence, proving the inequality.

Choose smooth metrics on the divisor line bundles and let $\beta_i$ denote their curvatures. The smooth closed form
\[
 \Theta=\pi^*(\beta|_Y)-\sum_{i=1}^q\tau_{Y,u}(E_i)\beta_i
\]
represents the displayed class. Poincar\'e--Lelong shows
$\int_Z\beta_i\wedge\gamma^{r-1}=\int_{E_i}\gamma^{r-1}$.
Thus $\int_Z\Theta\wedge\gamma^{r-1}\geq0$ for every Gauduchon metric on $Z$. For $r\geq2$, the generalized Lamari criterion of \cref{thm:lamari} supplies a genuine $w$ with $\Theta+\ddc w\geq0$. For $r=1$, choose a smooth probability form $\eta$ and solve the scalar Poisson equation $\ddc w=(\int_Z\Theta)\eta-\Theta$; the integral is nonnegative by the same inequality, so the resulting representative is positive. Since $\Theta$ is closed, this is a closed positive current, proving pseudoeffectivity.
\end{proof}

Equivalently, a strict violation of one of these finite Gauduchon inequalities forces $V_m=0$ for all sufficiently large $m$. Indeed the finite sum $\sum_i b_i(m)m^{-1}\int_{E_i}\gamma^{r-1}$ then exceeds the degree for every large $m$, contradicting the inequality satisfied by a nonzero restricted section.

The necessary condition has no general converse. On the elliptic curve $E=\C/(\mathbb Z+\mathrm{i}\mathbb Z)$, choose an irrational real number $a$ and the flat line bundle with multipliers $\mathrm{e}^{2\pi\mathrm{i} a}$ and $1$ on the two generators. Its constant metric has curvature zero. Take $u=0$ and $Y=X=E$, so every trace number is zero and the displayed pseudoeffectivity condition holds. Nevertheless every positive tensor power has no nonzero holomorphic section. To see this, a hypothetical section of a power has a zero divisor of degree zero and hence is nowhere zero. Its logarithmic norm is harmonic on $E$, and is constant. The lifted holomorphic function therefore has constant modulus and is constant on $\C$. The nontrivial multiplier $\mathrm{e}^{2\pi\mathrm{i} ma}$ forces that constant to be zero, a contradiction.
